\documentclass[11pt]{article}

\usepackage[margin=1in]{geometry}
\usepackage{amsmath,amssymb,mathtools}
\usepackage{array,booktabs,tabularx}
\usepackage{enumitem}
\usepackage{microtype}
\usepackage[colorlinks=true,linkcolor=blue,urlcolor=blue]{hyperref}

\newcommand{\Zp}{\mathbb{Z}_p}
\newcommand{\bits}{\{0,1\}}
\newcommand{\getsr}{\mathrel{\mathop{\leftarrow}\limits^{\$}}}
\newcommand{\negl}{\mathsf{negl}}
\newcommand{\PPT}{\mathsf{PPT}}
\newcommand{\inj}{\mathsf{inj}}
\newcommand{\lossy}{\mathsf{lossy}}
\newcommand{\Eval}{\mathsf{Eval}}
\newcommand{\KeyGen}{\mathsf{KeyGen}}
\newcommand{\Enum}{\mathsf{Enum}}
\newcommand{\Punc}{\mathsf{Punc}}
\newcommand{\Obf}{\mathcal{O}}
\newcommand{\Sign}{\mathsf{Sign}}
\newcommand{\Verify}{\mathsf{Verify}}
\newcommand{\Main}{\mathsf{Main}}
\newcommand{\Trig}{\mathsf{Trigger}}
\newcolumntype{Y}{>{\raggedright\arraybackslash}X}

\title{The Local Hybrid in the Salted and Derandomized\\Schnorr Instantiations}
\author{}
\date{}

\begin{document}
\maketitle

\begin{abstract}
This note explains the proof device that removes an a priori bound on the
number of signing queries in the salted and derandomized Schnorr
instantiations.  The proof changes one signing response at a time.  For that
one change it temporarily publishes a circuit containing two explicit hash
input--output pairs, exchanges an honest signature with a trigger signature,
and then restores the original circuit.  The temporary pairs do not remain
for the next query and therefore never accumulate.  A tunable lossy pre-hash
has a separate role: it lets the reduction guess the digest of an adaptive
message before the verification key is published.  Salt and derandomization
also have a separate role: they make repeated signing queries consistent with
the trigger simulation.  We give the salted proof first and then derive the
derandomized adaptive and static proofs from it.
\end{abstract}

\section{The point of the argument}

There are three different problems, and the construction uses a different
mechanism for each one.

The parameter \(r\) below is the upper bound on the lossy image of the
pre-hash.

\begin{enumerate}[leftmargin=2em]
\item A proof circuit cannot grow with the number of signing queries.  A
      \emph{local hybrid} treats one query at a time with a temporary
      fixed-size circuit.  This removes the query bound from the scheme.
\item An adaptive signing message is not known when the verification key is
      generated.  A tunable lossy function maps every message to one of at
      most \(r\) possible digests in the proof.  The reduction can therefore
      guess the digest at cost \(r\).
\item A trigger determines one simulated signature for each input at which
      it is evaluated.  Salt gives repeated invocations of the same message
      different trigger inputs.  Derandomization instead makes repeated
      invocations replay the same signature, so they may share one trigger
      input.
\end{enumerate}

The most important point is that there is no circuit that is gradually
punctured while one adversary runs.  A hybrid proof compares separate
experiments.  Each experiment publishes one fixed circuit at the start.  The
temporary circuit used to compare two adjacent experiments exists only in
that comparison.

\section{Common setup}

\subsection{Schnorr signatures}

Fix a security parameter \(\lambda\).  Let \(\mathbb{G}\) be a cyclic group
of prime order \(p\), let \(g\) be a generator, and let
\(\Zp=\{0,\ldots,p-1\}\) denote the exponents modulo \(p\).  The signing key is
\(x\getsr\Zp\), and the verification key contains
\[
  X=g^x.
\]
For a group element \(R\), a challenge \(c\in\Zp\), and a response
\(s\in\Zp\), the Schnorr verification equation is
\[
  g^s=R X^c.
\]
An honest signer chooses a nonce \(\rho\in\Zp\), sets \(R=g^\rho\), obtains a
challenge \(c\) from the hash function, and returns
\[
  s=\rho+cx \pmod p.
\]

\subsection{The tunable lossy pre-hash}

A tunable lossy function is a triple of algorithms
\[
  \mathsf{L}
  =
  \bigl(\mathsf{L.KeyGen},\mathsf{L.Eval},\mathsf{L.Enum}\bigr).
\]
The algorithms have the following meanings.

\begin{itemize}[leftmargin=2em]
\item
  \(\mathsf{L.KeyGen}(1^\lambda,r,\mu)\) takes an image parameter \(r\) and a
  mode \(\mu\in\{\inj,\lossy\}\).  It returns a public evaluation key \(K\)
  and private enumeration information \(\mathsf{td}\).
\item
  \(\mathsf{L.Eval}(K,m)\) evaluates the function on a message \(m\).
\item
  If \(K\) was generated in lossy mode, then
  \(\mathsf{L.Enum}(\mathsf{td})\) returns a list containing every possible
  output of \(\mathsf{L.Eval}(K,\cdot)\).  After duplicates are removed, the
  list has at most \(r\) entries.
\end{itemize}

Keys generated in injective mode and lossy mode are computationally
indistinguishable.  Except with a negligible key-generation failure
probability, evaluation in injective mode is one-to-one.  In the real scheme
the key is generated in injective mode.  The proof switches it to lossy mode.

Fix an injective encoding of the range of \(\mathsf{L.Eval}\) as bit strings
whose first bit is always 0.  A string whose first bit is 1 is therefore not a
valid digest encoding.
The \emph{digest} of a message \(m\) under key \(K\) is
\[
  d=\mathsf{dig}(K,m)
  :=
  \mathsf{encode}\bigl(\mathsf{L.Eval}(K,m)\bigr).
\]
The encoding is part of the definition of \(\mathsf{dig}\).  Thus
\(\mathsf{dig}\) is an algorithm, not a new assumption.

\subsection{The pseudorandom functions}

The construction uses four puncturable pseudorandom functions.  For
\(i\in\{1,2,3,4\}\), write
\[
  \mathsf{F}_i
  =
  \bigl(\mathsf{F}_i.\KeyGen,
        \mathsf{F}_i.\Eval,
        \mathsf{F}_i.\Punc\bigr).
\]
Algorithm \(\mathsf{F}_i.\KeyGen(1^\lambda)\) returns a key \(k_i\).
Algorithm \(\mathsf{F}_i.\Eval(k_i,z)\) returns the function value at \(z\).
For a finite set \(S\) of inputs,
\(\mathsf{F}_i.\Punc(k_i,S)\) returns a punctured key, denoted
\[
  k_i\{S\}.
\]
This key evaluates the same function at every input outside \(S\), but hides
the values at the inputs in \(S\).  Puncturable-PRF security lets the proof
replace those hidden values by independent uniform values.

The first two functions, \(\mathsf{F}_1\) and \(\mathsf{F}_2\), compute the
main branch of the hash.  The last two, \(\mathsf{F}_3\) and
\(\mathsf{F}_4\), compute one trigger challenge and response for each
\emph{slot}.  A slot will be a digest--salt pair in the salted scheme and a
digest or message in the derandomized schemes.  All four functions have range
\(\Zp\).

Fix two injective encoding algorithms:
\[
\begin{aligned}
  E_0(u)   &:\text{ encodes a slot \(u\) for }
              \mathsf{F}_3,\mathsf{F}_4,\\
  E_1(P,u) &:\text{ encodes a group element \(P\) and a slot \(u\) for }
              \mathsf{F}_1,\mathsf{F}_2.
\end{aligned}
\]
The encodings only put the indicated objects into the required PRF domains.
They have no cryptographic role beyond being injective.

The derandomized signer will also use an ordinary PRF
\[
  \mathsf{F}_0
  =
  \bigl(\mathsf{F}_0.\KeyGen,\mathsf{F}_0.\Eval\bigr)
\]
whose range is \(\Zp\).  This PRF need not be puncturable.

\subsection{Indistinguishability obfuscation}

Let \(\Obf\) be an indistinguishability obfuscator.  On input a circuit \(C\),
\(\Obf(1^\lambda,C)\) returns an obfuscated implementation of \(C\).
Security says that if two equal-size circuits compute exactly the same
function on every input, then their obfuscations are computationally
indistinguishable.  Every circuit below is padded to one fixed polynomial
size before it is obfuscated.

\subsection{The common hash circuit}

Fix an efficiently computable function
\[
  f:\mathbb{G}\longrightarrow\Zp
\]
such that \(f(Y)\ne 0\) for every \(Y\in\mathbb{G}\).  This is the conversion
function used in the circular discrete-logarithm problem.

For a public key \(X\), PRF keys \(k_1,\ldots,k_4\), a group element \(P\),
and a slot \(u\), define the circuit
\(\mathsf{C}[X,k_1,k_2,k_3,k_4]\) as follows:
\[
\begin{array}{ll}
t\gets E_0(u),&
c_{\mathrm{tr}}\gets\mathsf{F}_3.\Eval(k_3,t),\\
&
s_{\mathrm{tr}}\gets\mathsf{F}_4.\Eval(k_4,t),\\[2pt]
W\gets g^{s_{\mathrm{tr}}}X^{-c_{\mathrm{tr}}}.&
\end{array}
\]
If \(P=W\), the circuit returns \(c_{\mathrm{tr}}\).  Otherwise it sets
\[
\begin{aligned}
 e&=E_1(P,u),\\
 a&=\mathsf{F}_1.\Eval(k_1,e),\\
 b&=\mathsf{F}_2.\Eval(k_2,e),
\end{aligned}
\]
and returns
\[
  \bigl(f(PX^a g^b)+a\bigr)\bmod p.
\]
The first case is the trigger branch.  The second case is the main branch.

The trigger branch always contains a valid Schnorr signature.  Indeed,
\[
  W=g^{s_{\mathrm{tr}}}X^{-c_{\mathrm{tr}}}
  \quad\Longrightarrow\quad
  g^{s_{\mathrm{tr}}}=W X^{c_{\mathrm{tr}}}.
\]
The reduction can therefore return
\((W,s_{\mathrm{tr}})\) without knowing \(x\).

The main branch converts a forgery to a circular discrete-logarithm solution.
The circular discrete-logarithm problem gives the solver \(X\) and asks for
\((Y,z)\) satisfying
\[
  g^z=Y X^{f(Y)}.
\]
Suppose a valid forgery \((P,s)\) receives its challenge from the main branch.
The reduction knows the corresponding values \(a,b\), sets
\[
  Y=PX^a g^b,
\]
and returns \((Y,s+b)\).  If
\(c=f(PX^a g^b)+a=f(Y)+a\), then
\[
\begin{aligned}
  g^{s+b}
   &=P X^c g^b\\
   &=P X^{f(Y)+a}g^b\\
   &=Y X^{f(Y)}.
\end{aligned}
\]
Thus every valid forgery is useful: a main-branch forgery gives a circular
discrete-logarithm solution, while a trigger-branch forgery will give an
ordinary discrete logarithm after a separate puncturing step.

\section{Salted Schnorr}

\subsection{The scheme}

Fix a salt length \(\ell_{\mathrm{salt}}\).  Key generation samples
\(x\getsr\Zp\), sets \(X=g^x\), generates a tunable lossy-function key \(K\)
in injective mode, generates \(k_1,\ldots,k_4\), and publishes \(X\), \(K\),
and an obfuscation of the common hash circuit.

To sign a message \(m\), the signer samples independently
\[
  \rho\getsr\Zp
  \qquad\text{and}\qquad
  \sigma\getsr\bits^{\ell_{\mathrm{salt}}},
\]
and sets
\[
\begin{aligned}
  d&=\mathsf{dig}(K,m),&
  u&=(d,\sigma),&
  R&=g^\rho,\\
  c&=\mathsf{C}(R,u),&
  s&=\rho+cx\pmod p.
\end{aligned}
\]
The signature is \((R,s,\sigma)\).  Verification recomputes \(d\) and \(u\),
evaluates the published circuit to obtain \(c\), and checks
\(g^s=RX^c\).

\subsection{The outer experiments}

Fix a \(\PPT\) adversary that makes at most \(q_s\) signing queries.  Let
\(N\le q_s\) be the number of invocations made in one execution.  Number the
invocations that occur \(1,\ldots,N\), and let \(m_j\) be the message in
invocation \(j\).  The proof may sample \(q_s\) salts and honest nonces before key
generation:
\[
  \sigma_j\getsr\bits^{\ell_{\mathrm{salt}}}
  \qquad\text{and}\qquad
  \rho_j\getsr\Zp
  \qquad (j=1,\ldots,q_s).
\]
These values are independent coins.  Moving their sampling earlier does not
change their distribution, and unused values are discarded.

The proof also aborts on either of two events.  The first event occurs if two
signing invocations use the same slot.  The second occurs if two distinct
messages among the signing queries and forgery have the same digest.  In the
real experiment, an injective key rules out the second event, except on
key-generation failure, and the first event then requires a repeated salt.
The proof keeps these aborts after changing the key to lossy mode.  Therefore,
on every execution that continues, no other response uses the target slot
of query \(j\).

For \(i\in\{0,\ldots,q_s\}\), define \(G_i\) to be the complete experiment
in which the first \(i\) signing invocations that occur receive trigger
signatures and the remaining invocations receive honest signatures.  Every \(G_i\)
publishes an obfuscation of the original, unpunctured circuit.  The
experiments differ only in which signing procedure answers each invocation.

The proof compares \(G_0,G_1,\ldots,G_{q_s}\).  We now give the full
comparison between \(G_{j-1}\) and \(G_j\), restricted to executions in
which \(N\ge j\).  When \(N<j\), the two experiments are identical.

\subsection{Fixing the slot for query \texorpdfstring{\(j\)}{j}}

The digest and slot used by query \(j\) are
\[
  d_j=\mathsf{dig}(K,m_j)
  \qquad\text{and}\qquad
  u_j=(d_j,\sigma_j).
\]
The reduction knows \(\sigma_j\) before key generation, because it sampled
that independent coin itself.  It does not know \(m_j\), because the
adversary chooses it after seeing the verification key.  It therefore does
not know \(d_j\).

The proof first switches \(K\) from injective mode to lossy mode.  Using the
private information returned by lossy key generation, the reduction computes
\[
  D\gets\mathsf{L.Enum}(\mathsf{td}).
\]
It removes duplicates and pads the result with dummy strings to obtain a list
of exactly \(r\) entries.  The dummy strings are chosen outside the digest
encoding range.  The reduction selects a uniform list position, calls its
entry \(\widehat d_j\), and builds the query-\(j\) comparison for
\[
  \widehat u_j=(\widehat d_j,\sigma_j).
\]
It aborts unless \(\widehat d_j=d_j\).  Except with the enumeration failure
probability, \(d_j\) appears exactly once in the padded list, so
\[
  \Pr[\widehat d_j=d_j]=\frac1r.
\]
For the rest of this subsection, condition on a correct guess and write
\[
  u=u_j=\widehat u_j.
\]
This conditioning fixes the slot before the circuit is published.

\subsection{The honest and trigger records}

For the honest signature on query \(j\), define
\[
\begin{aligned}
  R&=g^{\rho_j},\\
  v&=\mathsf{C}(R,u),\\
  s_0&=\rho_j+vx\pmod p.
\end{aligned}
\]
The honest record is
\[
  T_0=(R,v,s_0).
\]

For the trigger signature at the same slot, define
\[
\begin{aligned}
  t&=E_0(u),\\
  c&=\mathsf{F}_3.\Eval(k_3,t),\\
  s_1&=\mathsf{F}_4.\Eval(k_4,t),\\
  W&=g^{s_1}X^{-c}.
\end{aligned}
\]
The trigger record is
\[
  T_1=(W,c,s_1).
\]
The actual response given to the adversary omits the challenge, since the
verifier recomputes it.  Thus the two possible responses are
\((R,s_0,\sigma_j)\) and \((W,s_1,\sigma_j)\).

The proof aborts if \(R=W\).  At either endpoint, \(R\) is uniform and
independent of \(W\), so this event has probability \(1/p\).  The two
endpoint aborts together cost at most \(2/p\).  From now on assume \(R\ne W\).

\subsection{The temporary two-point circuit}

Define the two main-branch inputs
\[
  e_R=E_1(R,u)
  \qquad\text{and}\qquad
  e_W=E_1(W,u),
\]
the puncture set
\[
  S=\{e_R,e_W\},
\]
and the two explicit hash input--output pairs
\[
  L=\{(e_R,v),(e_W,c)\}.
\]
The two entries of \(L\) are stored in a canonical order, for example in
lexicographic order of their first coordinates.  The order does not record
which entry came from the honest signature and which came from the trigger
signature.

The temporary circuit \(C_j\) contains
\[
  k_1\{S\},\quad k_2\{S\},\quad
  k_3\{\{t\}\},\quad k_4\{\{t\}\},
\]
together with \(t\) and \(L\).  On input a group element \(P\) and a slot
\(u'\), it proceeds as follows.

\begin{enumerate}[leftmargin=2em]
\item Compute \(e'=E_1(P,u')\).  If \((e',z)\in L\), return \(z\).
\item Compute \(t'=E_0(u')\).  If \(t'\ne t\), use the punctured keys
      \(k_3\{\{t\}\},k_4\{\{t\}\}\) to compute the trigger for \(u'\), and
      return its challenge if \(P\) is its trigger point.
\item Use \(k_1\{S\},k_2\{S\}\) to compute the main branch and return its
      output.
\end{enumerate}

The punctured keys are never evaluated at a punctured input.  In the first
step, the two missing complete hash outputs are supplied explicitly: \(v\)
is the main-branch output at \(R\), while \(c\) is the trigger-branch output
at \(W\).  The main-branch PRF outputs at \(W\) are neither supplied nor
needed.  When \(t'=t\), the circuit skips the trigger computation.  The
trigger point for this slot is \(W\), which the first step already handles.
Every other point at this slot fails the trigger comparison and should take
the main branch.

\subsection{Why the two circuits compute the same function}

At the true PRF values, the original circuit and \(C_j\) agree on every
input.

\begin{itemize}[leftmargin=2em]
\item On \((R,u)\), the original circuit takes the main branch because
      \(R\ne W\), and returns \(v\).  The temporary circuit returns the
      explicit value \(v\).
\item On \((W,u)\), the original circuit takes the trigger branch and returns
      \(c\).  The temporary circuit returns the explicit value \(c\).
\item On \((P,u)\) for \(P\notin\{R,W\}\), the original circuit fails its
      trigger comparison and takes the main branch.  The temporary circuit
      skips the unavailable trigger computation and takes the same main
      branch.
\item On every other slot, all punctured keys evaluate exactly as the
      original keys do.
\end{itemize}

The circuits are padded to the same size.  Indistinguishability obfuscation
therefore permits the proof to replace the published original circuit by an
obfuscation of \(C_j\).

\subsection{What each puncture accomplishes}

Puncturing and explicit storage have complementary roles.  Puncturing removes
the circuit's ability to recover a selected PRF output from the retained key.
Explicit storage preserves the hash function at the selected complete hash
inputs.  With the original values stored, the temporary circuit computes the
same function.  Once the relevant keys are punctured, puncturable-PRF security
lets the proof replace the hidden values by uniform ones.

The main-branch PRFs \(\mathsf{F}_1,\mathsf{F}_2\) are punctured at
\[
  e_R=E_1(R,u)
  \qquad\text{and}\qquad
  e_W=E_1(W,u).
\]
At \(e_R\), the two PRF outputs are the masks \(a,b\) used to produce the
honest challenge
\[
  v=f(RX^ag^b)+a\pmod p.
\]
Replacing \(b\) by a uniform exponent makes \(RX^ag^b\) uniform in the group
and independent of \(a\).  Replacing \(a\) by a uniform exponent then makes
\(v\) uniform in \(\Zp\), whatever the distribution of \(f\) on a uniform
group element.  Thus puncturing at \(e_R\) is what lets the proof turn the
honest record into a uniformly generated Schnorr record.

The main branch is not evaluated at \(W\), because the trigger branch returns
first.  Puncturing \(\mathsf{F}_1,\mathsf{F}_2\) at \(e_W\) is nevertheless
important.  If only \(e_R\) were punctured, the puncture set itself would
identify \(R\) as the honest point.  Puncturing at both \(e_R\) and \(e_W\)
makes the retained circuit description symmetric in the two point--value
pairs.  This is what allows the later exact exchange.  The values of the
main-branch PRFs at \(e_W\) are hidden for symmetry, not because the original
circuit uses them.

The trigger PRFs \(\mathsf{F}_3,\mathsf{F}_4\) are punctured at the one slot
input
\[
  t=E_0(u).
\]
Their outputs are the trigger challenge \(c\) and response \(s_1\).
Puncturing at \(t\) lets the proof replace both by independent uniform
values, so that
\[
  (W,c,s_1)=\bigl(g^{s_1}X^{-c},c,s_1\bigr)
\]
has the same form as the transformed honest record.  The temporary circuit
cannot evaluate the punctured trigger keys at \(t\), so it handles
\((W,u)\) through the explicit pair \((e_W,c)\) and skips the trigger
computation at every other point with slot \(u\).  This preserves the
function because \(W\) is the only point at that slot for which the trigger
comparison succeeds.

When a trigger output changes in its PRF hybrid, the proof also recomputes
\(W\), \(e_W\), the explicit pair containing \(W\), and the main-branch
puncture set.  This does not violate the selective puncturing requirement.
These are separate hybrid experiments: each one publishes one circuit with
its current values fixed at the start.  No already published circuit is
modified while the adversary runs.
The slot input \(t\) was fixed independently of the trigger-PRF keys.  After
the trigger challenge value determines \(W\), the corresponding reduction
generates the independent main-branch keys and punctures them at
\(\{e_R,e_W\}\).  Conversely, when a main-branch PRF is challenged, the set
\(\{e_R,e_W\}\) was produced without that PRF key.  Each puncture set is
therefore independent of the particular key whose security is being used,
even though not every point is fixed before every other key is generated.

In short, main-branch puncturing makes the honest record uniform and removes
its label; trigger-branch puncturing makes the trigger record uniform.  The
proof exchanges the records only after both jobs have been done.

\subsection{Making the two records symmetric}

Once the punctured circuit is published, puncturable-PRF security lets the
proof replace the selected \(\mathsf{F}_1\) and \(\mathsf{F}_2\) outputs at
\(e_R,e_W\) by independent uniform values.  Only the values at \(e_R\) are
used to form \(v\); both inputs are punctured and replaced so that the
circuit description treats \(R\) and \(W\) in the same way.

For fixed \(R\), if \(a,b\getsr\Zp\), then
\[
  U=RX^a g^b
\]
is uniform in \(\mathbb{G}\) and independent of \(a\).  Hence
\[
  v=f(U)+a\pmod p
\]
is uniform in \(\Zp\) and independent of \(R\).  The proof may replace \(v\)
by a fresh uniform value exactly, without a computational assumption.

Now change variables in the honest record.  Sampling
\(\rho_j,v\getsr\Zp\), setting \(R=g^{\rho_j}\), and setting
\(s_0=\rho_j+vx\) has exactly the same distribution as sampling
\(s_0,v\getsr\Zp\) and setting
\[
  R=g^{s_0}X^{-v}.
\]

Next, puncturable-PRF security at \(t\) replaces
\[
  c=\mathsf{F}_3.\Eval(k_3,t)
  \qquad\text{and}\qquad
  s_1=\mathsf{F}_4.\Eval(k_4,t)
\]
by independent uniform values.  The trigger point is recomputed as
\[
  W=g^{s_1}X^{-c}.
\]

At this midpoint the two records are
\[
  T_0=(g^{s_0}X^{-v},v,s_0)
  \qquad\text{and}\qquad
  T_1=(g^{s_1}X^{-c},c,s_1),
\]
Before conditioning, \(v,s_0,c,s_1\) are independent uniform elements of
\(\Zp\), so \(T_0\) and \(T_1\) are generated independently by the same
rule.  The event that their first components are different is invariant under
exchanging the two records.  Conditioning on that event therefore preserves
invariance of their joint distribution under exchange, although the records
need no longer be independent.

Before the exchange, the proof discards the original nonce \(\rho_j\), the
selected main-branch masks, the full PRF keys, and any unused copy of the two
responses.  It retains \(x\), the adversary's view before query \(j\), the
independent coins for every other signing query, the punctured keys, the
slot, and the canonically ordered pair set.  This state is enough to continue
the experiment.  Because of the repeated-slot abort, no other signing
response on an execution that continues uses \(u\).

The retained circuit contains the unordered set
\[
  \{(E_1(R,u),v),(E_1(W,u),c)\}
\]
and nothing that labels one entry ``honest'' and the other ``trigger.''  The
other retained state is unchanged if the two records are exchanged.  It is
therefore exact to return \((W,s_1,\sigma_j)\) instead of
\((R,s_0,\sigma_j)\).

\subsection{The complete query-\texorpdfstring{\(j\)}{j} comparison}

The adjacent-game comparison is
\[
\begin{array}{c}
G_{j-1}:
\text{ original circuit; query \(j\) receives the honest response}\\
\downarrow\quad \text{indistinguishability obfuscation}\\
\text{temporary circuit \(C_j\) with two explicit hash pairs}\\
\downarrow\quad \text{puncturable-PRF and exact distribution changes}\\
\text{the joint distribution of the two records is invariant under exchange}\\
\downarrow\quad \text{exchange the returned record}\\
\text{temporary circuit \(C_j\); query \(j\) receives the trigger response}\\
\downarrow\quad \text{reverse all preceding changes}\\
G_j:
\text{ original circuit; query \(j\) receives the trigger response}.
\end{array}
\]

The response is not exchanged immediately after the iO step.  It is
exchanged only at the symmetric midpoint, after the selected PRF values have
been hidden and the joint distribution of the records has become symmetric.
The
iO steps change the circuit description; the symmetry argument changes
which signature is returned.

\subsection{Why the query bound disappears}

Let \(Q\) denote the a priori signing-query bound that the earlier
construction placed in the definition of the scheme and in its circuit
padding.

After the proof reaches \(G_j\), the temporary circuit \(C_j\) has
disappeared.  The endpoint again publishes the original circuit.  To compare
\(G_j\) with \(G_{j+1}\), the proof starts over and constructs a new
temporary circuit \(C_{j+1}\) containing only the two hash pairs for query
\(j+1\).

Thus every circuit used in the proof contains at most two exceptional hash
inputs and one exceptional trigger slot.  These values never accumulate.
The circuit can be padded to one fixed polynomial size that does not depend
on the adversary's number of queries.

The proof still performs one adjacent-game comparison for each of the
adversary's \(q_s\) signing queries.  Its running time and security loss
therefore depend on \(q_s\).  This is ordinary: every polynomial-time
adversary has some polynomial query bound.  What has disappeared is a
parameter \(Q\) chosen when the scheme is defined and built into the circuit
size.  One fixed scheme now handles every polynomial-time adversary.

\subsection{What lossiness and salt do}

The local hybrid removes \(Q\), but it requires the slot \(u_j\) to be fixed
before the temporary circuit is published.  The reduction can presample the
salt \(\sigma_j\) and honest nonce \(\rho_j\).  It cannot presample the
adversary's message \(m_j\), so it guesses the digest \(d_j\).  Lossy mode
makes that guess cost \(r\) rather than the size of the message space.
Lossiness does not keep the circuit small; the local hybrid does.

Salt handles repeated messages.  If the adversary asks for the same message
twice, the two invocations receive independent salts.  Their slots are
\[
  (\mathsf{dig}(K,m),\sigma_j)
  \qquad\text{and}\qquad
  (\mathsf{dig}(K,m),\sigma_{j'}).
\]
Except with the salt-collision probability, the slots are different, so the
two trigger signatures are pseudorandomly fresh, matching the two fresh
honest randomized signatures.

\section{Derandomized Schnorr}

\subsection{The scheme and repeated messages}

The derandomized signer samples a nonce-PRF key
\[
  k_0\getsr\mathsf{F}_0.\KeyGen(1^\lambda)
\]
at key generation.  On message \(m\), it computes
\[
  \rho_m=\mathsf{F}_0.\Eval(k_0,m)
  \qquad\text{and}\qquad
  R_m=g^{\rho_m}.
\]
In the adaptive construction it also computes
\[
  d_m=\mathsf{dig}(K,m)
\]
and uses the digest \(d_m\) as the slot.  The published hash circuit is the
common circuit of Section~2.5 with \(u=d_m\).  The signer obtains
\[
  c_m=\mathsf{C}(R_m,d_m)
  \qquad\text{and}\qquad
  s_m=\rho_m+c_mx\pmod p,
\]
and returns \((R_m,s_m)\).  There is no salt in the signature.

The signing algorithm is deterministic.  A repeated query on \(m\) computes
the same nonce, point, challenge, and response.  The security game may
therefore answer it by replaying the first signature.  The proof indexes its
hybrid by the distinct queried messages, not by every invocation.

This is exactly what the trigger simulation needs.  For the slot \(d_m\),
the trigger PRFs determine one fixed challenge, response, and point:
\[
\begin{aligned}
  c^{\mathrm{tr}}_m
    &=\mathsf{F}_3.\Eval(k_3,E_0(d_m)),\\
  s^{\mathrm{tr}}_m
    &=\mathsf{F}_4.\Eval(k_4,E_0(d_m)),\\
  W_m
    &=g^{s^{\mathrm{tr}}_m}X^{-c^{\mathrm{tr}}_m}.
\end{aligned}
\]
Returning the same trigger signature on a repeated message now matches the
real signing behavior.

\subsection{The adaptive proof}

The proof first replaces the nonce-PRF outputs on the distinct queried
messages by independent uniform values.  This follows from ordinary PRF
security because \(k_0\) is used only for nonce derivation.  After this
change, every honest nonce point is uniform, exactly as in the salted proof.

Let \(N\le q_s\) be the number of distinct messages queried in one execution,
and let \(m_j\) be the \(j\)-th message when it is first queried.  For
\(i\in\{0,\ldots,q_s\}\), let \(G_i^{\mathrm{der}}\) be the experiment in
which the first \(i\) distinct messages that occur receive trigger signatures
and the remaining distinct messages receive honest signatures.  Repeats are
answered by replay.  Every experiment publishes the original unpunctured
circuit.  The comparison for index \(j\) is needed only when \(N\ge j\).

Fix the \(j\)-th distinct message, call it \(m_j\), and write
\[
  d_j=\mathsf{dig}(K,m_j).
\]
The reduction switches the pre-hash key to lossy mode, enumerates its image,
and guesses \(d_j\) before the verification key is published.  The guess
costs \(r\).  Conditioned on a correct guess, set
\[
\begin{aligned}
  u&=d_j,&
  t&=E_0(u),&
  \rho&=\rho_{m_j},&
  R&=g^\rho,\\
  c&=\mathsf{F}_3.\Eval(k_3,t),&
  s_1&=\mathsf{F}_4.\Eval(k_4,t),&
  W&=g^{s_1}X^{-c}.
\end{aligned}
\]
Next define
\[
\begin{aligned}
  e_R&=E_1(R,u),&
  e_W&=E_1(W,u),\\
  a&=\mathsf{F}_1.\Eval(k_1,e_R),&
  b&=\mathsf{F}_2.\Eval(k_2,e_R),\\
  v&=f(RX^ag^b)+a\pmod p,&
  s_0&=\rho+vx\pmod p.
\end{aligned}
\]
The temporary circuit contains the two explicit pairs
\[
  (e_R,v)
  \qquad\text{and}\qquad
  (e_W,c),
\]
punctures the first two PRFs at \(\{e_R,e_W\}\), and punctures the last two at
\(t\).  This is the same two-point circuit as in the salted proof, with the
digest \(d_j\) in place of the digest--salt pair \(u_j\).

The rest of the adjacent-game comparison is unchanged:

\begin{enumerate}[leftmargin=2em]
\item exclude the event \(R=W\);
\item use iO to publish the functionally equal two-point circuit;
\item use puncturable-PRF security and exact changes of variables to make
      the joint distribution of \((R,v,s_0)\) and \((W,c,s_1)\) invariant
      under exchange;
\item exchange which record is returned;
\item reverse the changes and restore the original circuit.
\end{enumerate}

The temporary circuit again disappears at the end of the step.  No punctures
or explicit pairs remain for the next distinct message.  Consequently the
derandomized adaptive scheme has no query bound in its circuit or key.

The two adaptive proofs differ only in how they handle repeated messages.
The salted proof assigns a different slot to each invocation.  The
derandomized proof assigns one slot to each distinct message and replays its
signature.

\subsection{The static proof}

In the static security game, the adversary declares all signing messages and
the forgery message before key generation.  The construction can omit the
pre-hash and use the message \(m\) itself as the slot.

For a declared signing message \(m_j\), the reduction already knows the slot
\[
  u=m_j
\]
before key generation.  It can also presample the uniform nonce point used
after the nonce-PRF hybrid.  The trigger point still depends on trigger-PRF
outputs, so it need not be fixed before every PRF key is generated.  What the
puncturable-PRF reductions require is that their challenge inputs be sampled
independently of the particular PRF key being challenged; the reductions
generate the other independent keys after receiving the challenged values.
No digest guess is needed.  The local hybrid still treats only one distinct
message at a time, so the exceptional pairs do not accumulate and the scheme
has no query parameter.

The static proof therefore has two useful properties at once:

\begin{itemize}[leftmargin=2em]
\item It removes the a priori query bound by the same local hybrid used in the
      adaptive proofs.
\item It needs no tunable lossy pre-hash because every slot is known before
      key generation.
\end{itemize}

This cleanly separates the two mechanisms.  The local hybrid is what removes
the scheme-level parameter \(Q\).  The lossy pre-hash is what lets the same
local hybrid be used when the messages are adaptive.

The proof still begins by replacing the deterministic nonce-PRF outputs on
the declared distinct messages by uniform nonces.  It then runs one local
hybrid per distinct message.  Since there is no factor \(r\) from guessing,
polynomially hard circular discrete logarithm, iO, and PRFs suffice.  The
number of declared messages affects the polynomial reduction loss, but it is
not a parameter of the scheme.

\section{Forgery branches}

After all signing responses have been changed to trigger responses, the
reduction no longer uses the secret exponent to answer signing queries.  It
can embed a hardness challenge in the public key and divide a forgery into
two cases.

\paragraph{Main branch.}
If the forgery is evaluated by the main branch, the algebra in Section~2.5
gives a circular discrete-logarithm solution.

\paragraph{Trigger branch.}
Suppose the forged point is the trigger point for its slot.  The reduction
first fixes the forged slot; in an adaptive proof it guesses the slot, while
in the static proof the declared forgery message determines it.  Call the
fixed slot \(u^\star\), and set
\[
  t^\star=E_0(u^\star)
  \qquad\text{and}\qquad
  c_{\mathrm{tr}}=\mathsf{F}_3.\Eval(k_3,t^\star).
\]
It punctures only the trigger-response PRF \(\mathsf{F}_4\) at \(t^\star\)
and uses iO to publish a functionally equal circuit carrying
\[
  Z=g^{\mathsf{F}_4.\Eval(k_4,t^\star)}
\]
explicitly.  At the guessed slot this circuit uses the trigger point
\(ZX^{-c_{\mathrm{tr}}}\); at every other slot it uses the original trigger
code.  With the true value of \(Z\), this computes the original function.

Puncturable-PRF security now replaces the hidden exponent at \(t^\star\) by
a uniform value, making \(Z\) uniform in \(\mathbb{G}\).  Only then does the
reduction set \(Z\) to an ordinary discrete-logarithm challenge \(Y=g^y\).
The trigger point becomes
\[
  W=Y X^{-c_{\mathrm{tr}}}.
\]
If the adversary returns a valid trigger-branch forgery with response \(s^*\),
then
\[
  g^{s^*}=W X^{c_{\mathrm{tr}}}=Y,
\]
so \(s^*=y\pmod p\) solves the discrete-logarithm challenge.

The challenge PRF \(\mathsf{F}_3\) is not punctured in this forgery
reduction.  Its output \(c_{\mathrm{tr}}\) remains known and is used to form
\(YX^{-c_{\mathrm{tr}}}\).  The discrete logarithm is embedded in the
response side of the trigger, which is why \(\mathsf{F}_4\) is the punctured
PRF here.  This is different from the signing swap, where both trigger PRFs
are punctured so that both \(c\) and \(s_1\) become uniform.

The cost of selecting the forged slot depends on the variant.

\begin{itemize}[leftmargin=2em]
\item In salted adaptive Schnorr, the forged slot contains both an adaptive
      digest and a salt chosen by the forger.  Guessing it costs
      \(r\cdot 2^{\ell_{\mathrm{salt}}}\).
\item In derandomized adaptive Schnorr, the forged slot is only the adaptive
      digest.  Guessing it costs \(r\).
\item In derandomized static Schnorr, the forgery message was declared before
      key generation.  Its slot is known, so there is no guessing cost.
\end{itemize}

\section{Digest collisions}

The real scheme uses an injective pre-hash key.  Except with the stated
key-generation failure probability, two distinct messages have different
digests.  Thus the pre-hash does not turn a signature on one message into a
signature on another.

The proof changes the key to lossy mode, where collisions must exist if the
message space is larger than the image.  The proof aborts if the adversary
finds two distinct relevant messages with the same digest.  This abort is
safe: an efficient algorithm that finds such a pair distinguishes lossy mode
from injective mode, since the pair cannot exist for a correctly generated
injective key.  No separate collision-resistant hash is being added.  The
needed collision guarantee follows from injectivity in the real mode and
indistinguishability of the two modes.

\section{Comparison}

\begin{center}
\small
\begin{tabularx}{\textwidth}{@{}>{\raggedright\arraybackslash}p{0.17\textwidth}
                              YYYY@{}}
\toprule
Result
& Signing
& Trigger slot
& Repeated message
& Why the proof knows the slot\\
\midrule
Salted adaptive
& Fresh random nonce and salt
& Digest and salt
& New slot, except on a salt collision
& Presample the salt and guess the digest from a lossy image of size \(r\)\\
\addlinespace
Derandomized adaptive
& Nonce PRF
& Digest
& Replay the same signature
& Guess the digest from a lossy image of size \(r\)\\
\addlinespace
Derandomized static
& Nonce PRF
& Message
& Replay the same signature
& All signing and forgery messages are declared before key generation\\
\bottomrule
\end{tabularx}
\end{center}

All three results use the same local two-point hybrid, so none has a
scheme-level signing-query bound.  The adaptive results use a lossy pre-hash
because the message is chosen after the verification key is published.  The
static result does not.  Salted signing preserves fresh randomized signatures
but adds a salt to the signature.  Derandomized signing preserves the
ordinary two-element Schnorr signature format but makes repeat queries
replay.

\section{Takeaway}

The proof is easiest to remember as a division of labor.

\begin{itemize}[leftmargin=2em]
\item The \textbf{temporary two-point circuit} changes one signing response
      without storing all signing responses.  This removes \(Q\).
\item The \textbf{tunable lossy pre-hash} turns one unknown adaptive message
      into one guess among \(r\) possible digests.  This handles adaptivity.
\item The \textbf{salt} gives each randomized signing invocation a new slot.
\item The \textbf{nonce PRF} makes a repeated message replay one signature,
      so one slot per distinct message is enough.
\end{itemize}

At the beginning and end of every adjacent-game comparison, the published
circuit is the original fixed-size circuit.  The proof never carries the
punctures for query \(j\) into the comparison for query \(j+1\).  That is the
precise sense in which the hybrid is local, and it is the reason the scheme
does not need to know an adversary's signing-query bound in advance.

\end{document}
